\subsection{Estados e Variáveis de Controle}
\label{sec:var_estado_controle}

Para este trabalho, adotou-se um conjunto de variáveis de estado, conforme:
\begin{equation}
    \label{eq:variaveis_estado}
var_{estado} = \{ \vec r,\; \vec v,\; q,\; \vec \omega,\; \theta,\; \dot\theta \},
\end{equation}
onde $\vec r$ e $\vec v$ representam posição e velocidade relativas no referencial LVLH, $q$ é o quaternion unitário da base, $\vec \omega$ a velocidade angular corporal, $\theta$ os ângulos articulares e $\dot\theta$ suas velocidades.
Por sua vez, as variáveis de controle utilizadas são:
\begin{equation}
    \label{eq:variaveis_controle}
F_{controle} = \{ \vec f,\; \vec{\tau}\},
\end{equation}
em que $\vec f$ são forças, $\vec{\tau}$ são os torques.
O modelo dinâmico inclui as equações orbitais, sendo o modelo do problema de dois corpos, as Equações de Hill/CW linearizadas, as equações de atitude do corpo rígido e as equações do manipulador robótico.

